% \DeclareDocumentCommand{\foocmd}{ O{default1} O{default2} m }{#1~#2~#3}
%                                     ⤷ #1        ⤷ #2     ⤷ #3
% usage: \foocmd[nondefault2][notfoo2]{foo}

% -------------------------------------------------- %
% Load the required packages if not loaded
% -------------------------------------------------- %
\makeatletter
\@ifpackageloaded{xparse}{}{\usepackage{xparse}}
\@ifpackageloaded{tabularx}{}{\usepackage{tabularx}} % for L R column type
\@ifpackageloaded{xifthen}{}{\usepackage{xifthen}}
\@ifpackageloaded{xspace}{}{\usepackage{xspace}} % for smart spaces
\@ifpackageloaded{amsopn}{}{\usepackage{amsopn}} % for \DeclareMathOperator
\@ifpackageloaded{amssymb}{}{\usepackage{amssymb}} % for \mathbb
\@ifpackageloaded{mathtools}{}{\usepackage{mathtools}} % for \ceil
\@ifpackageloaded{bm}{}{\usepackage{bm}} % for bold math
%\@ifpackageloaded{algorithm2e}{}{\usepackage[lined,linesnumbered,ruled]{algorithm2e}} % for pseudo code
\@ifpackageloaded{siunitx}{}{%
\usepackage[detect-all,scientific-notation=false,separate-uncertainty=true]{siunitx} % for proper handling of units
\sisetup{%
    retain-explicit-plus=true,% show Xe+2 instead of Xe2
    detect-weight=true,
    detect-family=true,
	output-exponent-marker=\ensuremath{\mathrm{e}},% usage: \num{1e-10}
	per-mode=symbol,% reciprocal (default) or symbol
	range-phrase={,}\ ,% Value intervals / ranges
    range-units=brackets,
    open-bracket={[},
    close-bracket=],
}%
}
\makeatother

% -------------------------------------------------- %
% Colors
% -------------------------------------------------- %
\definecolor{fMRTblack}{HTML}{2B2E34}
\definecolor{fMRTwhite}{HTML}{F3EEE1}
\definecolor{fMRTyellow}{HTML}{FFD564}
\definecolor{fMRTorange}{HTML}{FF5500}
\colorlet{fMRTlightGray}{white!80!fMRTblack}
\colorlet{fMRTgray}{white!35!fMRTblack}
\colorlet{fMRTdarkGray}{white!20!fMRTblack}
\colorlet{fMRTlightOrange}{fMRTorange!60!fMRTwhite}
\definecolor{fMRTblue}{HTML}{3333B3}
\colorlet{fMRTlightBlue}{fMRTblue!30!white}
\definecolor{fMRTverylightRed}{HTML}{FF8E6B}
\colorlet{fMRTlightRed}{red!60!fMRTwhite}
\colorlet{fMRTdarkRed}{red!70!fMRTgray} % red!60!fMRTgray
\definecolor{fMRTlightGreen}{HTML}{8CDD81} % green soap	
\definecolor{fMRTborwn}{HTML}{8B4513} % saddle brown
\colorlet{fMRTlightBorwn}{fMRTborwn!60!fMRTwhite}

% -------------------------------------------------- %
% Beamer Related
% -------------------------------------------------- %
% Background gradients
\colorlet{fMRTdark@backgroundInner}{fMRTblack}
\colorlet{fMRTdark@backgroundOuter}{fMRTblack}
\colorlet{fMRTlight@backgroundInner}{fMRTwhite}
\colorlet{fMRTlight@backgroundOuter}{fMRTwhite}

% Colorized emph commads
\newcommand{\emphred}[1]{\textcolor{fMRTlightRed}{#1}}
\newcommand{\emphgreen}[1]{\textcolor{fMRTlightGreen}{#1}}
\newcommand{\emphblue}[1]{\textcolor{fMRTblue}{#1}}
\newcommand{\emphyellow}[1]{\textcolor{fMRTyellow}{#1}}
\newcommand{\emphorange}[1]{\textcolor{fMRTorange}{#1}}
%\DeclareTextFontCommand{\emphred}{\itshape\color{fMRTdarkRed}} % https://tex.stackexchange.com/questions/47259/why-use-declaretextfontcommand-vs-just-newcommand

% Headings
\newcommand{\beamerheading}[1]{\par\bigskip{\Large\bfseries#1}\par\smallskip}

% -------------------------------------------------- %
% Tables
% -------------------------------------------------- %
% In combination with the tabularx environment
\newcolumntype{L}{>{\raggedright\arraybackslash}X}
\newcolumntype{R}{>{\raggedleft\arraybackslash}X}
\newcolumntype{C}{>{\centering\arraybackslash}X}
\newcolumntype{P}[1]{>{\raggedright\arraybackslash}p{#1}} % align left with fixed width
% For line breaks in cells (the optinal arg is t, c, or b to force the desired vertical alignment)
\newcommand{\breakeccell}[2][c]{\begin{tabular}[#1]{@{}c@{}}#2\end{tabular}}
\newcommand{\breakelcell}[2][c]{\begin{tabular}[#1]{@{}l@{}}#2\end{tabular}}
\newcommand{\breakercell}[2][c]{\begin{tabular}[#1]{@{}r@{}}#2\end{tabular}}

% -------------------------------------------------- %
% PGF
% -------------------------------------------------- %
%\newcommand\inputpgf[2]{{
%\let\pgfimageWithoutPath\pgfimage
%\renewcommand{\pgfimage}[2][]{\pgfimageWithoutPath[##1]{#1/##2}}
%\input{#1/#2}
%}}

% -------------------------------------------------- %
% Just for Editing
% -------------------------------------------------- %
% The following require the xcolor package; https://www.namsu.de/Extra/pakete/Xcolor.html
\newcommand{\predicted}[1]{\textcolor{ForestGreen}{#1}}
\newcommand{\alter}[1]{\textcolor{Blue}{#1}}
\newcommand{\unsure}[1]{\textcolor{BrickRed}{#1}}
\newcommand{\refneeded}[1]{\textcolor{BrickRed}{\textbf{[add ref: #1]}}}

% -------------------------------------------------- %
% Flags
% -------------------------------------------------- %
\newcommand{\Next}{^\prime} % e.g. next state
\newcommand{\optsym}{\star} % symbol to indicate optimality
\newcommand{\opt}{^\optsym}
\newcommand{\obssym}{\text{obs}}
\newcommand{\obs}{^\obssym} % observations in the context of LFI
\newcommand{\initsym}{\text{init}}
\newcommand{\init}{_\initsym}
\newcommand{\simusym}{\text{sim}}
\newcommand{\simu}{^\simusym}
\newcommand{\realsym}{\text{real}}
\newcommand{\real}{^\realsym}

% -------------------------------------------------- %
% Math Related
% -------------------------------------------------- %
\newcommand{\RR}{\mathbb{R}} % the set of real numbers to the power of #1
\newcommand{\onehalf}{\frac{1}{2}}
\newcommand{\onethird}{\frac{1}{3}}
\newcommand{\onefourth}{\frac{1}{4}}
\newcommand{\alphad}{\dot{\alpha}} % not bold
\newcommand{\betad}{\dot{\beta}} % not bold
% A set is a collection of items that is unordered and does not allow duplicates
\DeclareDocumentCommand{\set}{O{} O{} m}{\{#3\}_{#1}^{#2}} % #1 = subscript, #2 = superscipt, #3 = content
% A tuple is an ordered sequence of items that allows duplicates
\DeclareDocumentCommand{\tuple}{O{} O{} m}{\left\langle#3\right\rangle_{#1}^{#2}} %  #1 = subscript, #2 = superscipt, #3 = content
\newcommand{\data}{{\cal D}} % data
\newcommand{\order}[1]{\mathcal{O}\!\left( #1\right) }

% MDP --------------------------------------------- %
\newcommand{\model}{\mathcal{M}}%
\newcommand{\stateset}[1][]{\mathcal{S}_{#1}} % set of states in a MDP; Note: using \stateset as a subscript requires curly brackets, i.e. _{\stateset}
\newcommand{\obsset}[1][]{\mathcal{O}_{#1}} % set of observations in a MDP
\newcommand{\actionset}[1][]{\mathcal{A}_{#1}} % set of actions in a MDP
\newcommand{\transprobsym}{\mathcal{P}}
\makeatletter
\DeclareDocumentCommand{\@transprob}{O{} O{} m}{\transprobsym_{#1}^{#2}\left( #3\right) }  % starred version
\DeclareDocumentCommand{\@@transprob}{O{} O{} m}{\transprobsym_{#1}^{#2}\!\left( #3\right) }
\newcommand{\transprob}{\@ifstar\@transprob\@@transprob}
\makeatother
\newcommand{\initstatedistrsym}{\mu}
\makeatletter
\DeclareDocumentCommand{\@initstatedistr}{O{} O{} m}{\initstatedistrsym_{#1}^{#2}\left( #3\right) }  % starred version
\DeclareDocumentCommand{\@@initstatedistr}{O{} O{} m}{\initstatedistrsym_{#1}^{#2}\!\left( #3\right) }
\newcommand{\initstatedistr}{\@ifstar\@initstatedistr\@@initstatedistr}
\makeatother
\newcommand{\mdp}[1][]{\mathcal{M}_{#1}} % initial state distribution in MDP
\newcommand{\dimstate}{{n_s}} % d
\newcommand{\dimobs}{{n_o}} % q
\newcommand{\dimact}{{n_a}} % p
%\newcommand{\dimmdlparam}{m}
\newcommand{\dimpolparam}{{n_{\polparam}}} % n
\newcommand{\dimdomparam}{{n_{\domparam}}} % m

\newcommand{\estopt}{_{\text{est}}\opt} % optimizer of the etimator
\newcommand{\trueopt}{_{\text{true}}\opt} % optimmzer of the true function
\newcommand{\candsym}{{n_c}}
\newcommand{\refsym}{{n_r}}
\newcommand{\Hinfty}{\mathcal{H}_\infty} % H infinity robust control

% -------------------------------------------------- %
% Abbreviations for Text
% -------------------------------------------------- %
\newcommand{\etal}{et al.~}
\newcommand{\eg}{e.g.\xspace}
\newcommand{\ie}{i.e.\xspace}
\newcommand{\cf}{cf.\xspace}
\newcommand{\iid}{i.i.d.\xspace}
\newcommand{\st}{s.t.\xspace}
\newcommand{\wrt}{w.r.t.\xspace}
\newcommand{\aka}{a.k.a.\xspace}
\newcommand{\suppmat}{Supplementary Material}
\newcommand{\SimtoSim}{{Sim-to-Sim}\xspace}
\newcommand{\Simtosim}{{Sim-to-sim}\xspace}
\newcommand{\simtosim}{{sim-to-sim}\xspace}
\newcommand{\SimtoReal}{{Sim-to-Real}\xspace}
\newcommand{\Simtoreal}{{Sim-to-real}\xspace}
\newcommand{\simtoreal}{{sim-to-real}\xspace}
\newcommand{\sota}{{state-of-the-art}\xspace}
\newcommand{\Sota}{{State-of-the-art}\xspace}

\newcommand{\ballbalancersym}{\textrm{bb}}
\newcommand{\cartpolesym}{\textrm{cp}}
\newcommand{\futuapendulumsym}{\textrm{fp}}
\newcommand{\sub}{swing-up and balance\xspace}
\newcommand{\Sub}{Swing-up and balance\xspace}
\newcommand{\bic}{ball-in-a-cup\xspace}
\newcommand{\Bic}{Ball-in-a-cup\xspace}

% -------------------------------------------------- %
% Symbols for Operators
% -------------------------------------------------- %
% https://tex.stackexchange.com/questions/67506/newcommand-vs-declaremathoperator
% \DeclareMathOperator can only be used in the preamble while \newcommand has no such restriction.

\DeclareMathOperator*{\argmax}{arg\,max}
\DeclareMathOperator*{\argmin}{arg\,min}

\newcommand{\distr}[3]{#1\left( #2 \big| #3\right) } % probability distribution: #1 symbol, #2 random variable, 3# parameters
\newcommand{\distrnormal}[2]{\distr{\mathcal{N}}{#1}{#2}} % normal distribution: #1 random variable, #2 parameters, e.g. mean, covar
\newcommand{\distruniform}[2]{\distr{\mathcal{U}}{#1}{#2}} % uniform distribution: #1 random variable, #2 parameters, e.g. min, max
\newcommand{\distrbinom}[2]{\distr{\mathcal{B}}{#1}{#2}} % binomial distribution: #1 random variable, #2 parameters, e.g. n, p
\newcommand{\distrbern}[2]{\distr{\mathcal{B}}{#1}{#2}} % Bernoulli distribution: #1 random variable, #2 parameters, e.g. p
\newcommand{\p}[1]{p\left( #1\right) } % probability
\newcommand{\E}[1]{\mathbb{E}\! \left[ #1\right] } % expected value of a random variable or a function of a random variable
\newcommand{\Esub}[2]{\mathbb{E}_{#1} \! \left[ #2\right] } % expected value with subscript
\newcommand{\Esubbig}[2]{\mathbb{E}_{#1} \! \big[ #2\big] } % expected value with subscript (enforce big brackets)
\newcommand{\EsubBig}[2]{\mathbb{E}_{#1} \! \Big[ #2\Big] } % expected value with subscript (enforce Big brackets)
\newcommand{\V}[1]{\mathbb{V} \! \left[ #1\right] } % variance of a random variable or a function of a random variable
\newcommand{\Vsub}[2]{\mathbb{V}_{#1} \! \left[ #2\right] } % variance with subscript
\newcommand{\Cov}[1]{\mathrm{Cov}\left( #1\right) } % covariance between two random variables or two functions of random variables
\renewcommand{\Pr}[1]{\mathbb{P}\left( {#1}\right) } % probability of an event

\makeatletter
\newcommand{\@KL}[2]{\mathrm{KL}(#1 \Vert #2)} % starred version; #1 and #2 are probability distribution
\newcommand{\@@KL}[3][]{\mathrm{KL}\!\left( #2 #1\Vert #3\right)} % #1(optional) = size of dilimiter, 1 and #2 are probability distribution e.g. \KL*{\small}{a}{b}
\newcommand{\KL}{\@ifstar\@KL\@@KL}
\makeatother
\newcommand{\KLlosssym}{\mathcal{L}_{\textrm{KL}}}
\newcommand{\quantilesym}{Q} % or \mathrm{Q} or \mathbb{Q}
\newcommand{\quantile}[2]{\quantilesym_{#1} \! \left[ #2\right] } % quantile of a set: #1 level, #2 set
\newcommand{\bias}[1]{\mathrm{b}[#1] } % bias | former \left[#1 \right]
\newcommand{\logb}[2]{\mathrm{log}_{#1}(#2)} % logarithm to a specific base (#1)
\newcommand{\entropy}[1]{\mathrm{H}\!\left( #1\right) }
\renewcommand{\ln}[1]{\mathrm{ln}\left( #1\right) } % logarithm
\renewcommand{\log}[1]{\mathrm{log}\left( #1\right) } % logarithm
\renewcommand{\exp}[1]{\mathrm{exp}\left( #1\right) } % exponential function
\newcommand{\fracpartial}[2]{\frac{\partial #1}{\partial #2}}
%\newcommand{\diff}[1]{\ensuremath{\operatorname{d}\!{#1}}}
\newcommand{\diff}{\mathrm{d}}
\renewcommand{\sin}[1]{\mathrm{sin}\!\left( #1\right) }
\newcommand{\sinsq}[1]{\mathrm{sin}^2\!\left( #1\right) }
\renewcommand{\cos}[1]{\mathrm{cos}\!\left( #1\right) }
\newcommand{\cossq}[1]{\mathrm{cos}^2\!\left( #1\right) }
%https://tex.stackexchange.com/questions/448/how-to-automatically-resize-the-vertical-bar-in-a-set-comprehension
\makeatletter
\newcommand{\@given}[2]{\left. #1\right| #2} % starred version; #1 = random variable, #2 = condition
\newcommand{\@@given}[3][]{#2 #1| #3} % #1(optional) = size of dilimiter, #2 = random variable, #3 = condition, e.g. \given*{\small}{a}{b}
\newcommand{\given}{\@ifstar\@given\@@given}
\makeatother

% Measurements
\newcommand{\mean}[1]{\mu\!\left( #1\right)}
\newcommand{\estmean}[1]{m\!\left( #1\right)}
\newcommand{\std}[1]{\sigma\!\left( #1\right)}
\newcommand{\eststd}[1]{s\!\left( #1\right)}
\newcommand{\halfspan}[1]{h\!\left( #1\right)}

\DeclareDocumentCommand{\distmeas}{O{} O{} m}{\text{D}_{#1}^{#2} \!\left( #3\right)} % 2 optional args

% Just for convenience
\newcommand{\agivens}{\fa | \fs}
\newcommand{\ninit}{n_{\text{init}}}
\newcommand{\tholdsucc}{\edrsym^{\text{succ}}}
\newcommand{\tholdsuccQQ}{375}

\DeclarePairedDelimiter{\ceil}{\lceil}{\rceil} % larger integer
\DeclarePairedDelimiter{\floor}{\lfloor}{\rfloor} % samller integer

% -------------------------------------------------- %
% Matrix Related
% -------------------------------------------------- %
\newcommand{\matele}[3]{\left[ #1 \right]_{#2#3}} % element of a matrix
\newcommand{\matwsz}[3]{\underset{(#2 \times #3)}{#1}} % matrix with size beneath
\newcommand{\tran}{^\textsf{\upshape T}} % transposed
\newcommand{\inv}{^{-1}} % inverse
\newcommand{\invtran}{^{-\textsf{\upshape T}}} % inverse and transposed
\newcommand{\pinv}{^{+}} % pseudo-inverse
\newcommand{\rank}[1]{\mathrm{rank}\left( #1\right) } % rank operator, perhaps \mathup{rank} is better
%%%\renewcommand{\det}[1]{\left| {#1} \right|} % determinant
\newcommand{\trace}[1]{\operatorname{tr}\left( #1\right) }
\newcommand{\adj}[1]{\mathrm{adj}\left( #1\right) } % adjungate
\newcommand{\diag}[1]{\mathrm{diag}\left( #1\right) } % diagonal matrix

% Reward Function ---------------------------------- %
\newcommand{\rewsym}{r}
\DeclareDocumentCommand{\rewfcn}{O{} O{} m}{\rewsym_{#1}^{#2}\!\left( #3\right) }

% Trajectory --------------------------------------- %
\newcommand{\trajsym}{\tau}
\newcommand{\trajssym}{\bm{\tau}}
\newcommand{\trajspace}{\mathcal{T}}
\newcommand{\trajsetsym}{\mathcal{T}}
\DeclareDocumentCommand{\traj}{O{} O{}}{\trajssym_{#1}^{#2}} % 2 optional args
\newcommand{\statedistrsym}{\mu}
\DeclareDocumentCommand{\statedistr}{O{} O{} m}{\statedistrsym_{#1}^{#2}\!\left( #3\right) } % steady state distribution

% Domain Parameters -------------------------------- %
\newcommand{\domparam}{\xi}
\newcommand{\domparams}{\bm{\xi}}
\newcommand{\domparamspace}{\Xi}
\newcommand{\domparamsnom}{\bar{\bm{\xi}}}
\newcommand{\domparamdistrsym}{p} % prob dist for the physics params
\makeatletter
\DeclareDocumentCommand{\@domparamdistr}{O{} O{} m}{\domparamdistrsym_{#1}^{#2}\left( #3\right) }  % starred version
\DeclareDocumentCommand{\@@domparamdistr}{O{} O{} m}{\domparamdistrsym_{#1}^{#2}\!\left( #3\right) }
\newcommand{\domparamdistr}{\@ifstar\@domparamdistr\@@domparamdistr}
\makeatother

% Domain Parameter Distributions ------------------- %
\newcommand{\domdistrparam}{\phi}
\newcommand{\domdistrparams}{\bm{\phi}} % parameters of the distribution of the domain params
\newcommand{\domdistrparamspace}{\Phi}
\newcommand{\densestsym}{q}
\DeclareDocumentCommand{\densest}{O{} O{} m}{\densestsym_{#1}^{#2}\!\left( #3\right)} % 2 optional args; usage: \densest[1][2]{3}
\newcommand{\densestparam}{\phi}
\newcommand{\densestparams}{\bm{\phi}} % parameters of the density estimator of the domain params
\newcommand{\densestparamspace}{\Phi}
\DeclareDocumentCommand{\domparamprior}{O{} O{} m}{\domparamdistrsym_{#1}^{#2}\!\left( #3\right) } % e.g. p(xi)
\DeclareDocumentCommand{\domparamproposal}{O{} O{} m}{\tilde{\domparamdistrsym}_{#1}^{#2}\!\left( #3\right) } % e.g. p_tilde(xi)
\DeclareDocumentCommand{\domparamlikelihood}{O{} O{} m}{\domparamdistrsym_{#1}^{#2}\!\left( #3\right) } % e.g. p(x | xi)
\DeclareDocumentCommand{\domparamposterior}{O{} O{} m}{\domparamdistrsym_{#1}^{#2}\!\left( #3\right) } % e.g. p(xi | x)
\DeclareDocumentCommand{\estdomparamposterior}{O{} O{} m}{\hat{\domparamdistrsym}_{#1}^{#2}\!\left( #3\right) } % e.g. p_hat(xi | x)

% SBI related -------------------------------------- %
\newcommand{\genmodelsym}{g}
\DeclareDocumentCommand{\genmodel}{O{} O{} m}{\genmodelsym{#1}^{#2}\!\left( #3\right) }
\newcommand{\sbidata}{\bm{x}}
\newcommand{\sbidataspace}{\mathcal{X}}
\newcommand{\summstatsym}{f}
\newcommand{\summstatparam}{\psi}
\newcommand{\summstatparams}{\bm{\psi}}
\newcommand{\summstatparamspace}{\Psi}
\DeclareDocumentCommand{\summstatfcn}{O{} O{} m}{\summstatsym_{#1}^{#2}\!\left( #3\right) }

% Policy and its Parameters ------------------------ %
\newcommand{\polsym}{\pi}
\DeclareDocumentCommand{\pol}{O{} O{} m}{\polsym_{#1}^{#2}\!\left( #3\right) } % policy; optional arg for dependency on physics params
\newcommand{\polparam}{\theta}
\newcommand{\polparamspace}{\Theta}
\DeclareDocumentCommand{\polparams}{O{} O{}}{\ftheta_{#1}^{#2}} % 2 optional args; NOTE: max_\polparams yields an error, but max_{\polparams} is fine
\newcommand{\solcandi}{\polparams[][c]} % solution candidate
\DeclareDocumentCommand{\polparamsopt}{O{} O{}}{\ftheta_{#1}^{#2 \optsym}} % optimal policy parameters indicated by star superscript
\DeclareDocumentCommand{\polparamopt}{O{} O{}}{\polparam_{#1}^{#2 \optsym}} % optional arg is the index (not perfect - it would be better to define a new \DeclareDocumentCommand)
%\DeclareDocumentCommand{\polparamsopt}{O{} O{}}{\polparams[#1][#2 \optsym]} % optional arg is the index (not perfect - it would be better to define a new \DeclareDocumentCommand)
\newcommand{\polparamsoptcand}{\polparams[\candsym][\optsym]} % opt cand sol
\newcommand{\polparamsoptref}{\polparams[\refsym][\optsym]} % opt ref sol

% Gaussian Process --------------------------------- %
\newcommand{\GPsym}{\mathcal{GP}}
\newcommand{\GPmean}{m}
\newcommand{\GPcovar}{k}
\newcommand{\GP}[2]{\GPsym\!\left( #1,#2\right)}
\newcommand{\acqfcnsym}{a}
\DeclareDocumentCommand{\acqfcn}{O{} O{} m}{\acqfcnsym_{#1}^{#2}\!\left( #3\right) }
\newcommand{\boinputspace}[1][]{\mathcal{X}_{#1}} % input space for Bayesian optimization

% Optimality Gap ----------------------------------- %
\newcommand{\optgapsym}{G}
\DeclareDocumentCommand{\optgap}{O{} O{} m}{\optgapsym_{#1}^{#2}\!\left( #3\right) } % optimality gap of the policy #3; e.g. \optgap[1][2]{3}
\newcommand{\estoptgapsym}{\hat{G}}
\DeclareDocumentCommand{\estoptgap}{O{} O{} m}{\estoptgapsym_{#1}^{#2}\!\left( #3\right) } %  estimator of the optimality gap of the policy #3 calculated from #1 samples; e.g. \optgap[1][2]{3}
\newcommand{\meanoptgapsym}{\bar{G}}
\DeclareDocumentCommand{\optgapmean}{O{} O{} m}{\meanoptgapsym_{#1}^{#2}\!\left( #3\right) } % sample mean of the optimality gap from n_G calculations of G_n
\newcommand{\optgapsampsetsym}{\mathcal{G}}
\DeclareDocumentCommand{\optgapsampset}{O{} O{}}{\optgapsampsetsym_{#1}^{#2}} % set of optimality gap samples

% Simulation Optimization Bias --------------------- %
\newcommand{\sob}[1]{\mathrm{b}\!\left[#1 \right]}
\newcommand{\sobBig}[1]{\mathrm{b}\!\Big[#1 \Big]}
\newcommand{\sobNormal}[1]{\mathrm{b} [#1]}

% Lagrange function -------------------------------- %
\newcommand{\lagrangefcn}[1]{\mathcal{L}\left( #1\right) }

% Bootstrapped optimality gap ---------------------- %
\newcommand{\numbssamples}{{n_b}}
\iffalse
% old notation
\newcommand{\bootsym}{*}
\DeclareDocumentCommand{\bootestoptgap}{O{} O{} m}{{}^{\bootsym}\estoptgapsym_{#1}^{#2}\!\left( #3\right) }
\DeclareDocumentCommand{\bootmeanoptgap}{O{} O{} m}{{}^{\bootsym}\meanoptgapsym_{#1}^{#2}\!\left( #3\right) }
\DeclareDocumentCommand{\bootoptgapsampset}{O{} O{}}{{}^{\bootsym}\optgapsampsetsym_{#1}^{#2}}
\fi
\iftrue
% new notation
\newcommand{\bootsym}{B}
\DeclareDocumentCommand{\bootestoptgap}{O{} O{} m}{\estoptgapsym_{#1}^{\bootsym #2}\!\left( #3\right) }
\DeclareDocumentCommand{\bootmeanoptgap}{O{} O{} m}{\meanoptgapsym_{#1}^{\bootsym #2}\!\left( #3\right) }
\DeclareDocumentCommand{\bootoptgapsampset}{O{} O{}}{\optgapsampsetsym_{#1}^{\bootsym #2}}
\fi

% Return ------------------------------------------- %
\newcommand{\edrsym}{J}
\DeclareDocumentCommand{\edr}{O{} O{} m}{\edrsym_{#1}^{#2}\!\left( #3\right)} % expected discounted return
\newcommand{\optedrval}{j\opt}
\newcommand{\estedrsym}{\hat{J}}
\DeclareDocumentCommand{\estedr}{O{}  O{} m}{\estedrsym_{#1}^{#2}\!\left( #3\right)} % estimator of the expected (trajs) discounted return
\newcommand{\eedrsym}{J} % former \bar{J}
\DeclareDocumentCommand{\eedr}{O{} O{} m}{\eedrsym_{#1}^{#2}\!\left( #3\right)} % expected (domains) expected (trajs) discounted return
\DeclareDocumentCommand{\esteedr}{O{} O{} m}{\estedrsym_{#1}^{#2}\!\left( #3\right)} % estimator of the expected (domains) expected (trajs) discounted return
\newcommand{\objfunsym}{f}
\DeclareDocumentCommand{\objfun}{O{} m}{\objfunsym_{#1}\!\left( #2\right)} % objective function
\newcommand{\optfunval}{z\opt}

% Dynamical Systems & Potentials ------------------- %
\newcommand{\nummps}{n_{M}} % \numds = \nummps
\newcommand{\dsdynparams}{\fpsi}
\newcommand{\dsdynsym}{f}
\DeclareDocumentCommand{\dsdyn}{O{} m}{\dsdynsym_{#1}\!\left( #2\right)} % 1 optional arg (subscript)
\newcommand{\acti}{a}
\newcommand{\actis}{\fa}
\newcommand{\pot}{p}
\newcommand{\pots}{\fp}
\newcommand{\potdynparam}{\chi}
\newcommand{\potdynparams}{\fchi}
\newcommand{\potdynsym}{f}
\newcommand{\potsdynsym}{\ff}
\DeclareDocumentCommand{\potdyn}{O{} m}{\potdynsym_{#1}\!\left( #2\right)} % 1 optional arg (subscript)
\DeclareDocumentCommand{\potsdyn}{O{} m}{\potsdynsym_{#1}\!\left( #2\right)} % 1 optional arg (subscript)
\newcommand{\stim}{s}
\newcommand{\stims}{\fs}
\newcommand{\capacity}{C}
\newcommand{\sigmoid}[1]{\sigma\!\left( #1\right)}

% Robot related ------------------------------------ %
\newcommand{\jointpos}{q}
\newcommand{\jointvel}{\dot{q}}
\newcommand{\jointspos}{\bm{q}}
\newcommand{\jointsvel}{\dot{\jointposs}}
\newcommand{\taskpos}{x}
\newcommand{\taskvel}{\dot{\taskpos}}
\newcommand{\taskposs}{\bm{x}}
\newcommand{\taskvels}{\dot{\taskposs}}
\newcommand{\tasksposs}{\bm{X}}
\newcommand{\tasksvels}{\dot{\tasksposs}}

% Reference multiple equations --------------------- %
\ExplSyntaxOn
\NewDocumentCommand{\eqsref}{m}{\quinn_eqsref:n {#1}}
\seq_new:N \l_quinn_eqsref_seq
\cs_new:Npn \quinn_eqsref:n #1
{
	\seq_set_split:Nnn \l_quinn_eqsref_seq { , } { #1 }
	\seq_pop_right:NN \l_quinn_eqsref_seq \l_tmpa_tl
	( % print the left parenthesis
	\seq_map_inline:Nn \l_quinn_eqsref_seq
	{ \ref{##1},\nobreakspace } % print the first references
	\exp_args:NV \ref \l_tmpa_tl % print the last or only one
	) % print the right parenthesis
}
\ExplSyntaxOff

% -------------------------------------------------- %
% Pseudocode
% -------------------------------------------------- %
% see https://tex.stackexchange.com/questions/229355/algorithm-algorithmic-algorithmicx-algorithm2e-algpseudocode-confused
% and https://tex.stackexchange.com/questions/174631/using-algorithm2e-package-in-beamer
% https://tex.stackexchange.com/questions/200730/how-to-have-the-algorithmicx-comment-symbol-in-algorithm2e/200736
\SetKw{kwNot}{not}
\SetKw{kwAnd}{and}
\SetKw{kwBreak}{break}
\SetKwInOut{Input}{input}
\SetKwInOut{Output}{output}
\SetKwInOut{Initialization}{init}
\SetKwComment{Comment}{$\triangleright$\ }{}
\newcommand\famuracommfont[1]{\rmfamily\textcolor{darkgray}{#1}}
\SetCommentSty{famuracommfont}
\SetKwRepeat{Do}{do}{while} % do ... while ... loop
\SetKwBlock{With}{with}{end} % usage: \With(condition){code}; see p.12 in https://mlg.ulb.ac.be/files/algorithm2e.pdf
% inline if/else blocks without newline at the end
\newcommand{\llIf}[2]{{\let\par\relax\lIf{#1}{#2}}}
\newcommand{\llElse}[1]{{\let\par\relax\lElse{#1}}}
\newcommand{\llFor}[2]{{\let\par\relax\lFor{#1}{#2}}}

% Pseudocode: Subroutines -------------------------- %
\newcommand{\PolOptsym}{\textsc{PolOpt}}
\newcommand{\DistrOptsym}{\textsc{DistrOpt}}
\newcommand{\Infersym}{\textsc{Infer}}
\newcommand{\PolOpt}[1]{\PolOptsym\!\left[ #1\right]}
\newcommand{\DistrOpt}[1]{\DistrOptsym\!\left[ #1\right]}
\newcommand{\Infer}[1]{\Infersym\!\left[ #1\right]}

% -------------------------------------------------- %
% Bold Math Symbols
% -------------------------------------------------- %
% Latin Symbols ------------------------------------ %
\DeclareBoldMathCommand{\fnull}{0}
\DeclareBoldMathCommand{\fO}{0}
\DeclareBoldMathCommand{\fA}{A}
\DeclareBoldMathCommand{\fa}{a}
\DeclareBoldMathCommand{\fB}{B}
\DeclareBoldMathCommand{\fb}{b}
\DeclareBoldMathCommand{\fC}{C}
\DeclareBoldMathCommand{\fc}{c}
\DeclareBoldMathCommand{\fD}{D}
\DeclareBoldMathCommand{\fd}{d}
\DeclareBoldMathCommand{\fE}{E}
\DeclareBoldMathCommand{\fe}{e}
\DeclareBoldMathCommand{\fF}{F}
\DeclareBoldMathCommand{\ff}{f}
\DeclareBoldMathCommand{\fG}{G}
\DeclareBoldMathCommand{\fg}{g}
\DeclareBoldMathCommand{\fH}{H}
\DeclareBoldMathCommand{\fh}{h}
\DeclareBoldMathCommand{\fI}{I}
\DeclareBoldMathCommand{\fJ}{J}
\DeclareBoldMathCommand{\fK}{K}
\DeclareBoldMathCommand{\fM}{M}
\DeclareBoldMathCommand{\fm}{m}
\DeclareBoldMathCommand{\fN}{N}
\DeclareBoldMathCommand{\fn}{n}
\DeclareBoldMathCommand{\fo}{o}
\DeclareBoldMathCommand{\fP}{P}
\DeclareBoldMathCommand{\fp}{p}
\DeclareBoldMathCommand{\fQ}{Q}
\DeclareBoldMathCommand{\fq}{q}
\DeclareBoldMathCommand{\fq}{q}
\newcommand{\fqd}{\dot{\fq}}
\newcommand{\fqdd}{\ddot{\fq}}
\DeclareBoldMathCommand{\fR}{R}
\DeclareBoldMathCommand{\fr}{r}
\newcommand{\frd}{\dot{\fr}}
\DeclareBoldMathCommand{\fS}{S}
\DeclareBoldMathCommand{\fs}{s}
\newcommand{\fstilde}{\tilde{\fs}}
\DeclareBoldMathCommand{\ft}{t}
\DeclareBoldMathCommand{\fT}{T}
\DeclareBoldMathCommand{\fU}{U}
\DeclareBoldMathCommand{\fu}{u}
\DeclareBoldMathCommand{\fV}{V}
\DeclareBoldMathCommand{\fv}{v}
\DeclareBoldMathCommand{\fW}{W}
\DeclareBoldMathCommand{\fw}{w}
\DeclareBoldMathCommand{\fx}{x}
\DeclareBoldMathCommand{\fz}{z}
\newcommand{\fxd}{\dot{\fx}}
\newcommand{\fxdd}{\ddot{\fx}}
\newcommand{\fxbar}{\bar{\fx}}
\newcommand{\fxbard}{\dot{\fxbar}}
\newcommand{\fxbardd}{\ddot{\fxbar}}
\DeclareBoldMathCommand{\fY}{Y}
\DeclareBoldMathCommand{\fy}{y}
\DeclareBoldMathCommand{\fZ}{Z}
% Greek Symbols ------------------------------------ %
\DeclareBoldMathCommand{\falpha}{\alpha}
\newcommand{\falphad}{\dot{\falpha}}
\DeclareBoldMathCommand{\fbeta}{\beta}
\DeclareBoldMathCommand{\fchi}{\chi}
\DeclareBoldMathCommand{\fepsilon}{\epsilon}
\DeclareBoldMathCommand{\fvarepsilon}{\varepsilon}
\DeclareBoldMathCommand{\fgamma}{\gamma}
\DeclareBoldMathCommand{\fGamma}{\Gamma}
\DeclareBoldMathCommand{\flambda}{\lambda}
\DeclareBoldMathCommand{\fLambda}{\Lambda}
\DeclareBoldMathCommand{\fmu}{\mu}
\DeclareBoldMathCommand{\fnu}{\nu}
\DeclareBoldMathCommand{\fomega}{\omega}
\DeclareBoldMathCommand{\fpi}{\pi}
\DeclareBoldMathCommand{\fphi}{\phi}
\DeclareBoldMathCommand{\fPhi}{\Phi}
\DeclareBoldMathCommand{\fpsi}{\psi}
\DeclareBoldMathCommand{\fsigma}{\sigma}
\DeclareBoldMathCommand{\fSigma}{\Sigma}
\DeclareBoldMathCommand{\ftau}{\tau}
\DeclareBoldMathCommand{\ftheta}{\theta}
\DeclareBoldMathCommand{\fTheta}{\Theta}
\newcommand{\fhtheta}{\hat{\ftheta}}
\newcommand{\fthetad}{\dot{\ftheta}}
\newcommand{\fthetadd}{\ddot{\ftheta}}
\DeclareBoldMathCommand{\fxi}{\xi}
